\section{Déterminant angulaire et relation électrofaible de Weinberg}

Le déterminant de \(T\) est

\begin{equation}
\det T=q_+q_-=D_A.
\label{eq:detT}
\end{equation}

Le secteur constitutif du vide évalue \(F(T)^2\), où
\(F(z)=(1-z^{90})/(1-z^{120})\), tandis que le secteur angulaire
électrofaible évalue \(\det T\). Il s'agit de deux fonctions d'un
antécédent commun, et non d'une factorisation horizontale entre constantes.
Cette distinction est développée dans le \cref{ch:ew}.

\begin{remark}[Contrôle négatif]
Le bloc est réfuté si l'une quelconque de ces identités exactes est violée :
positivité de \(I-T^{120}\), égalité
\cref{eq:three-four-completion}, relations \cref{eq:vacuum-relations} ou
compatibilité \(E_m(\mathcal V_{m+1})=\mathcal V_m\). Une coïncidence SI
ne peut pas compenser une défaillance opératorielle.
\end{remark}

\subsection*{Portée logique et domaine de validité}
Opérateur, spectre, complétion \(3\to4\), parité, amplification et sections constitutives : \emph{Théorème interne exact}. La formulation qui les réunit en un seul opérateur est une \emph{Formalisation nouvelle} ; les données angulaires et les droites sont un \emph{Résultat antécédent démontré} du développement HMT précédent et sont déclarées ici comme structure de départ explicite.

